\documentclass[12pt,a4paper]{article}

\usepackage[a4paper,margin=1in]{geometry}
\usepackage{enumitem}
\usepackage{listings}
\usepackage{xcolor}
\usepackage{titlesec}

\definecolor{codegreen}{RGB}{0,120,0}
\definecolor{codegray}{RGB}{100,100,100}

\lstset{
    basicstyle=\ttfamily\small,
    keywordstyle=\color{blue},
    commentstyle=\color{codegreen},
    numberstyle=\tiny\color{codegray},
    numbers=left,
    frame=single,
    breaklines=true,
    showstringspaces=false
}

\title{\textbf{Operating Systems Laboratory}\\
Week--1 Practice Sheet\\
\Large Demonstration of \texttt{fork()} System Call}
\author{}
\date{}

\begin{document}

\maketitle

\section*{Instructions}

\begin{itemize}
    \item Write all programs in C under a Linux environment.
    \item Compile each program using:
\begin{lstlisting}
gcc filename.c -o output
./output
\end{lstlisting}
    \item Predict the output before execution.
    \item Draw the process tree wherever applicable.
    \item Write suitable comments in every program.
\end{itemize}

%%%%%%%%%%%%%%%%%%%%%%%%%%%%%%%%%%%%%%%%%%%%%%%%%%%%%%%

\section*{Classwork}

\begin{enumerate}

\item Write a program using \texttt{fork()} to create a single child process. Display:
\begin{itemize}
    \item Process ID (PID)
    \item Parent Process ID (PPID)
    \item Return value of \texttt{fork()}
\end{itemize}

\item Write a program that prints:
\begin{itemize}
    \item ``I am Parent Process''
    \item ``I am Child Process''
\end{itemize}
using the return value of \texttt{fork()}.

\item Write a program that creates one child process and displays the process hierarchy using \texttt{getpid()} and \texttt{getppid()}.

\item Execute the following program and predict the output before running.

\begin{lstlisting}
fork();
printf("Operating System Lab\n");
\end{lstlisting}

Explain why the output is printed twice.

\item Write a program containing:

\begin{lstlisting}
fork();
fork();
\end{lstlisting}

Predict:
\begin{itemize}
    \item Total number of processes created.
    \item Number of output lines.
    \item Draw the corresponding process tree.
\end{itemize}

\item Execute the following program and verify the output.

\begin{lstlisting}
fork();
fork();
fork();
printf("Linux\n");
\end{lstlisting}

Determine:
\begin{itemize}
    \item Number of processes.
    \item Number of times ``Linux'' is printed.
\end{itemize}

\item Write a program in which:
\begin{itemize}
    \item Child prints numbers from 1 to 5.
    \item Parent waits for the child using \texttt{wait()}.
    \item Parent prints ``Child Completed''.
\end{itemize}

\item Modify the previous program by removing \texttt{wait()}. Observe and explain the change in output.

\item Write a program that creates three child processes. Each child should print its PID and the parent should display the total number of children created.

\item Accept an integer from the user. The child process computes its square while the parent displays the original number.

\end{enumerate}

%%%%%%%%%%%%%%%%%%%%%%%%%%%%%%%%%%%%%%%%%%%%%%%%%%%%%%%

\section*{Homework}

\begin{enumerate}[resume]

\item Write a program where the parent prints all even numbers from 1 to 20 and the child prints all odd numbers.

\item Accept a number from the user. The child computes its factorial while the parent waits for the child to finish.

\item Accept an array of integers. The child computes the sum of all elements while the parent computes the average.

\item Write a program in which the child prints all prime numbers between 1 and 100 while the parent prints all composite numbers.

\item Write a program that generates the Fibonacci sequence using the child process.

\item Predict the output without executing.

\begin{lstlisting}
printf("A\n");
fork();
printf("B\n");
fork();
printf("C\n");
\end{lstlisting}

State the number of times A, B and C are printed.

\item Determine the total number of processes created by the following program. Draw the process tree.

\begin{lstlisting}
fork();

if(fork())
    fork();
\end{lstlisting}

\item Predict the number of outputs produced by the following program.

\begin{lstlisting}
fork();

if(fork())
{
    fork();
}

printf("OS Lab\n");
\end{lstlisting}

\item Write a program that recursively creates exactly ten processes. Each process should display:
\begin{itemize}
    \item PID
    \item PPID
    \item Level in the process tree
\end{itemize}

\item Create exactly fifteen processes using the minimum possible number of \texttt{fork()} calls.

\item Create a complete binary process tree of height four using \texttt{fork()}.

\item Determine the total number of processes created after the following program executes.

\begin{lstlisting}
fork();
fork();
fork();
fork();
\end{lstlisting}

Also derive the general formula for the number of processes created after \(n\) successive \texttt{fork()} calls.

\end{enumerate}

%%%%%%%%%%%%%%%%%%%%%%%%%%%%%%%%%%%%%%%%%%%%%%%%%%%%%%%

\section*{Discussion Questions}

\begin{enumerate}

\item Why does the \texttt{fork()} system call return twice?

\item What are the possible return values of \texttt{fork()}?

\item Explain the difference between the parent and child process.

\item What happens if \texttt{fork()} fails?

\item What is the purpose of the \texttt{wait()} system call?

\item Explain the concepts of orphan and zombie processes.

\item What is Copy-on-Write (CoW) and why is it used in \texttt{fork()}?

\item Differentiate between \texttt{fork()} and \texttt{exec()}.

\item Why is \texttt{getppid()} useful?

\item Explain how the operating system schedules the parent and child processes after a \texttt{fork()} call.

\end{enumerate}

\end{document}